\begin{table}[H]
    \centering
    \caption{Diagnostik Deviasi $SD$ Tersirat dari Korelasi Item-Total terhadap $SD$ Nasional per Subtes}
    \label{tab:diagnostik}
    \begin{spacing}{1.0}
    \begin{tabular}{lccccc}
        \toprule
        Subtes & $SD$ Sampel & $r_{iz}$ & $SD$ Tersirat Populasi & $SD$ Nasional & Rasio \\
        \midrule
        PU              & 76,85 & 0,706 & 87,27 & 71,52 & 1,220 \\
        PPU             & 61,50 & 0,590 & 67,45 & 100,00 & 0,674 \\
        PBM             & 76,62 & 0,631 & 85,03 & 100,00 & 0,850 \\
        PK              & 131,90 & 0,705 & 149,75 & 100,00 & 1,498 \\
        LBI Saintek     & 82,31 & 0,586 & 90,15 & 100,00 & 0,902 \\
        LBI Soshum      & 70,11 & 0,452 & 74,15 & 100,00 & 0,742 \\
        LBE             & 86,50 & 0,642 & 96,32 & 100,00 & 0,963 \\
        PM              & 100,01 & 0,634 & 111,10 & 100,00 & 1,111 \\
        \bottomrule
    \end{tabular}
    \end{spacing}
    \vspace{2mm}
    \begin{flushleft}
    \footnotesize\textit{Catatan.} Rasio = $SD$ tersirat populasi $\div$ $SD$ nasional. Rasio yang menyimpang jauh dari 1 mengindikasikan pelanggaran asumsi korelasi seragam yang disyaratkan koreksi Case III untuk subtes itu.
    \end{flushleft}
\end{table}
